\documentclass[11pt]{article}
\usepackage{fontspec,amsmath,amssymb,amsthm,hyperref}
\usepackage[margin=1in]{geometry}
\setlength{\emergencystretch}{3em}
\title{MYPGD pointwise recursion and its exact Markov law}
\author{LangevinLean mathematical exposition}
\date{}
\begin{document}\maketitle
\noindent Source module: \texttt{ProximalLangevin/MYPGDAlgorithm.lean}. Compile this proof document with XeLaTeX or LuaLaTeX. The accompanying HTML pairs the mathematical argument with the exact Lean statements and records the scope of the formalization.


MYPGD is the expanded MYIPLA drift with only latent Gaussian noise retained. The source definitions use actual pairwise proximity and independent smoothing and time parameters.


\section{Mask only the innovations in the pointwise chain}


Define the update by replacing \(\xi\) in MYIPLA with \(P_x\xi\). The zero and successor recursion identities show inductively that the resulting MYPGD chain equals the MYIPLA pointwise chain driven by \((P_x\xi_k)_k\), and equals the generic noise-driven chain for its own update. Its law is the pushforward of this actual chain.




\textbf{Correspondence with the displayed declarations.}


\begin{itemize}

\item \texttt{mypgdUpdate}: MYPGD removes precisely the parameter noise from the displayed MYIPLA step.
\item \texttt{mypgdChain}: The actual MYPGD chain driven by supplied full innovations.
\item \texttt{mypgdChain\_eq\_myiplaChain}: MYPGD is exactly MYIPLA applied to the latent-only innovations.
\item \texttt{mypgdLaws}: The output law of the actual MYPGD recursion.
\end{itemize}


\noindent Lean declarations: \texttt{mypgdUpdate}, \texttt{mypgdChain}, \texttt{mypgdChain\_eq\_myiplaChain}, \texttt{mypgdChain\_eq\_noiseDrivenChain}, \texttt{mypgdLaws}.

\section{Recognize the shared deterministic map}


Let \(F_\eta=\sum_i(g_1+g_2^\eta)(\theta,x_i)\) in scaled coordinates and \(h=\gamma/N\). The displayed MYPGD step is exactly


\[
X_{k+1}=X_k-h\nabla F_\eta(X_k)+\sqrt{2h}\,P_x\xi_k.
\]


Its common parameter coordinate has no Gaussian perturbation; each physical latent coordinate has the usual \(\sqrt{2\gamma}\) noise. The expanded proximal-residual formula is identified with this gradient step. The pointwise chain is then identified with its measurable Gaussian Markov recursion.



Initialization has finite second moment, and the standard Gaussian vectors are independent across coordinates and times and independent of initialization. The parameter coordinates of those vectors are erased in MYPGD. They remain available as auxiliary noise for the comparison with MYIPLA. No initial iid-particle hypothesis is needed for this weaker finite-particle theorem.



The auxiliary measurability estimate uses the unconditional bound \(\operatorname{Lip}(I-(\gamma/N)\nabla F_\eta)\le1+|\gamma/N|\,NL_\eta\). It needs no contraction restriction and supplies measurable updates even before the convergence regime is imposed.




\textbf{Correspondence with the displayed declarations.}


\begin{itemize}

\item \texttt{mypgdUpdate\_eq\_latentGaussianUpdate}: The drift/noise formula for the actual Moreau proximal construction.
\item \texttt{myiplaUpdate\_eq\_fullGaussianUpdate}: The full-noise update used for comparison is the actual MYIPLA update.
\end{itemize}


\noindent Lean declarations: \texttt{mypgdUpdate\_eq\_latentGaussianUpdate}, \texttt{myiplaUpdate\_eq\_fullGaussianUpdate}, \texttt{smoothed\_forward\_lipschitz}.

\section{Independent coordinate Gaussians give the Markov recursion}


The Gaussian-innovation hypotheses imply that vector innovations have standard Gaussian law, are independent over time, and are independent of initialization. The measurable-chain law theorem therefore identifies the two pushforwards with the actual iterated Markov transitions \(D+aP_x\xi\) and \(D+a\xi\). This transfers coupling estimates from the canonical product space to any admissible algorithm realization.




The two update identities recognize the same deterministic map \(D(z)=z-(\gamma/N)\nabla V_{N,\eta}(z)\). MYPGD adds \(\sqrt{2\gamma/N}P_x\xi\), while MYIPLA adds the full \(\sqrt{2\gamma/N}\xi\). The coordinate-Gaussian assumptions give the complete independent standard vector-noise path and independence from initialization. The shared iid-chain law theorem consequently identifies each supplied implementation with its iterated Markov transition. This equality is needed before the abstract Markov comparison can be cited for these algorithms.




\textbf{Correspondence with the displayed declarations.}


\begin{itemize}

\item \texttt{mypgdLaws\_eq\_markovLaws}: The actual Gaussian realization has the proved latent-noise Markov law.
\item \texttt{myiplaLaws\_eq\_markovLaws}: The same full innovations give the ordinary full-noise MYIPLA Markov law.
\end{itemize}


\noindent Lean declarations: \texttt{mypgdLaws\_eq\_markovLaws}, \texttt{myiplaLaws\_eq\_markovLaws}.

\end{document}
